% GENERATED FILE. Edit canonical lessons, then run npm run generate:books.
% canonical-source-hash: fnv1a64:95d776c6b56283c0
% canonical-lessons: SA-C447-kedarah, SA-C447-sasyam, SA-C447-urvara, SA-C447-lunati, SA-C447-avacinoti, SA-R447-review-441-447, SA-R447-review-final-447

\chapter{Fields on the Way Home}
\label{ch:sa-fields-on-the-way-home}
\hlchaptermodality{447}

\begin{chapteropening}
\textbf{By the end of this chapter:} \emph{I can describe the fields and the harvest on the road home.}

Complete the last lesson of chapter 447: I can describe the fields and the harvest on the road home.
\end{chapteropening}

\section[\texorpdfstring{\sk{केदारः} (kedāraḥ)}{केदारः (kedāraḥ)}]{\sk{केदारः} (kedāraḥ) — a flooded field; a paddy field}

\label{lesson:SA-C447-kedarah}



\begin{quote}
Before the new one: say the Sanskrit for kusha grass, then the Sanskrit for a young shoot.
\end{quote}

\subsection*{You'll want to know: \sk{केदारः}}
\textbf{\sk{केदारः}} — \emph{kedāraḥ} — \textquotedblleft{}a flooded field; a paddy field\textquotedblright{}.

\begin{quote}
\sk{केदारेषु} \sk{जलम्} \sk{अस्ति।}
\end{quote}

\begin{quote}
\emph{kedāreṣu jalam asti.}
\end{quote}

\begin{quote}
There is water in the flooded fields.
\end{quote}

A masculine noun.

\begin{cousinweb}
\textbf{\sk{केदारः}} (\emph{kedāraḥ}) is the word MW gives for a field or meadow, especially one under water; MW gives no derivation, so its origin is uncertain (MW).
\end{cousinweb}

\begin{grammarlens}[title={What you've built}]
One more word for getting around.
\end{grammarlens}

\subsection*{Your turn}
\begin{itemize}
\raggedright
  \item \emph{Say it:} \emph{kedāraḥ}
  \item \emph{Say it:} \emph{kedāraḥ}, once more
  \item \emph{Say it:} \emph{kedāreṣu jalam asti.} — There is water in the flooded fields.
  \item \emph{Recall:} say the Sanskrit for kusha grass, then the Sanskrit for a young shoot, then say \emph{kedāraḥ} again
\end{itemize}

\begin{tcolorbox}[breakable,colback=teal!4,colframe=teal!35!black,title={Before you move on}]
What does \sk{केदारः} mean? (\textquotedblleft{}A flooded field; a paddy field\textquotedblright{}.) Say it once more.
\end{tcolorbox}
\section[\texorpdfstring{\sk{सस्यम्} (sasyam)}{सस्यम् (sasyam)}]{\sk{सस्यम्} (sasyam) — a crop}

\label{lesson:SA-C447-sasyam}



\begin{quote}
Before the new one: say the Sanskrit for a young shoot, then the Sanskrit for a flooded field.
\end{quote}

\subsection*{You'll want to know: \sk{सस्यम्}}
\textbf{\sk{सस्यम्}} — \emph{sasyam} — \textquotedblleft{}a crop\textquotedblright{}.

\begin{quote}
\sk{केदारेषु} \sk{हरितं} \sk{सस्यं} \sk{वर्धते।}
\end{quote}

\begin{quote}
\emph{kedāreṣu haritaṃ sasyaṃ vardhate.}
\end{quote}

\begin{quote}
A green crop grows in the flooded fields.
\end{quote}

A neuter noun.

\begin{cousinweb}
\textbf{\sk{सस्यम्}} (\emph{sasyam}) is the word MW gives for corn, grain, the fruit of the field, a crop; MW gives no derivation, so its origin is uncertain (MW; CDIAL 13295).
\end{cousinweb}

\begin{grammarlens}[title={What you've built}]
One more word for getting around.
\end{grammarlens}

\subsection*{Your turn}
\begin{itemize}
\raggedright
  \item \emph{Say it:} \emph{sasyam}
  \item \emph{Say it:} \emph{sasyam}, once more
  \item \emph{Say it:} \emph{kedāreṣu haritaṃ sasyaṃ vardhate.} — A green crop grows in the flooded fields.
  \item \emph{Recall:} say the Sanskrit for a young shoot, then the Sanskrit for a flooded field, then say \emph{sasyam} again
\end{itemize}

\begin{tcolorbox}[breakable,colback=teal!4,colframe=teal!35!black,title={Before you move on}]
What does \sk{सस्यम्} mean? (\textquotedblleft{}A crop\textquotedblright{}.) Say it once more.
\end{tcolorbox}
\section[\texorpdfstring{\sk{उर्वरा} (urvarā)}{उर्वरा (urvarā)}]{\sk{उर्वरा} (urvarā) — fertile land}

\label{lesson:SA-C447-urvara}



\begin{quote}
Before the new one: say the Sanskrit for a flooded field, then the Sanskrit for a crop.
\end{quote}

\subsection*{You'll want to know: \sk{उर्वरा}}
\textbf{\sk{उर्वरा}} — \emph{urvarā} — \textquotedblleft{}fertile land\textquotedblright{}.

\begin{quote}
\sk{कृषकः} \sk{उर्वरायां} \sk{धान्यं} \sk{वपति।}
\end{quote}

\begin{quote}
\emph{kṛṣakaḥ urvarāyāṃ dhānyaṃ vapati.}
\end{quote}

\begin{quote}
The farmer sows grain in the fertile land.
\end{quote}

A feminine noun.

\begin{cousinweb}
\textbf{\sk{उर्वरा}} (\emph{urvarā}) is probably connected, as MW suggests, with \emph{uru}, \textquotedblleft{}wide\textquotedblright{}, so its origin is uncertain; MW gives \emph{urvarā} as fertile soil, a field yielding crops (MW).
\end{cousinweb}

\begin{grammarlens}[title={What you've built}]
One more word for getting around.
\end{grammarlens}

\subsection*{Your turn}
\begin{itemize}
\raggedright
  \item \emph{Say it:} \emph{urvarā}
  \item \emph{Say it:} \emph{urvarā}, once more
  \item \emph{Say it:} \emph{kṛṣakaḥ urvarāyāṃ dhānyaṃ vapati.} — The farmer sows grain in the fertile land.
  \item \emph{Recall:} say the Sanskrit for a flooded field, then the Sanskrit for a crop, then say \emph{urvarā} again
\end{itemize}

\begin{tcolorbox}[breakable,colback=teal!4,colframe=teal!35!black,title={Before you move on}]
What does \sk{उर्वरा} mean? (\textquotedblleft{}Fertile land\textquotedblright{}.) Say it once more.
\end{tcolorbox}
\section[\texorpdfstring{\sk{लुनाति} (lunāti)}{लुनाति (lunāti)}]{\sk{लुनाति} (lunāti) — reaps a crop; cuts}

\label{lesson:SA-C447-lunati}



\begin{quote}
Before the new one: say the Sanskrit for a crop, then the Sanskrit for fertile land.
\end{quote}

\subsection*{You'll want to know: \sk{लुनाति}}
\textbf{\sk{लुनाति}} — \emph{lunāti} — \textquotedblleft{}reaps a crop; cuts\textquotedblright{}.

\begin{quote}
\sk{कृषकः} \sk{दात्रेण} \sk{सस्यं} \sk{लुनाति।}
\end{quote}

\begin{quote}
\emph{kṛṣakaḥ dātreṇa sasyaṃ lunāti.}
\end{quote}

\begin{quote}
The farmer reaps the crop with a sickle.
\end{quote}

Past (he …): \sk{सः} \sk{सस्यं} \sk{लूनवान्}, \emph{saḥ sasyaṃ lūnavān}.

\begin{cousinweb}
\textbf{\sk{लुनाति}} (\emph{lunāti}) comes from the root √lū, \textquotedblleft{}to cut\textquotedblright{}; MW gives \emph{lunāti} as to cut, sever, reap, with the participle \emph{lūna}, \textquotedblleft{}reaped\textquotedblright{}; MW compares Greek \emph{luō}, \textquotedblleft{}to loose\textquotedblright{} (MW).
\end{cousinweb}

\begin{grammarlens}[title={What you've built}]
One more word for getting around.
\end{grammarlens}

\subsection*{Your turn}
\begin{itemize}
\raggedright
  \item \emph{Say it:} \emph{lunāti}
  \item \emph{Say it:} \emph{lunāti}, once more
  \item \emph{Say it:} \emph{kṛṣakaḥ dātreṇa sasyaṃ lunāti.} — The farmer reaps the crop with a sickle.
  \item \emph{Recall:} say the Sanskrit for a crop, then the Sanskrit for fertile land, then say \emph{lunāti} again
\end{itemize}

\begin{tcolorbox}[breakable,colback=teal!4,colframe=teal!35!black,title={Before you move on}]
What does \sk{लुनाति} mean? (\textquotedblleft{}Reaps a crop; cuts\textquotedblright{}.) Say it once more.
\end{tcolorbox}
\section[\texorpdfstring{\sk{अवचिनोति} (avacinoti)}{अवचिनोति (avacinoti)}]{\sk{अवचिनोति} (avacinoti) — picks fruit; plucks flowers}

\label{lesson:SA-C447-avacinoti}



\begin{quote}
Before the new one: say the Sanskrit for fertile land, then the Sanskrit for reaps a crop.
\end{quote}

\subsection*{You'll want to know: \sk{अवचिनोति}}
\textbf{\sk{अवचिनोति}} — \emph{avacinoti} — \textquotedblleft{}picks fruit; plucks flowers\textquotedblright{}.

\begin{quote}
\sk{बालिका} \sk{वृक्षात्} \sk{फलानि} \sk{अवचिनोति।}
\end{quote}

\begin{quote}
\emph{bālikā vṛkṣāt phalāni avacinoti.}
\end{quote}

\begin{quote}
The girl picks fruit from the tree.
\end{quote}

Past (she …): \sk{सा} \sk{फलानि} \sk{अवचितवती}, \emph{sā phalāni avacitavatī}.

\begin{cousinweb}
\textbf{\sk{अवचिनोति}} (\emph{avacinoti}) is built from \emph{ava}, \textquotedblleft{}down\textquotedblright{}, and the root √ci, \textquotedblleft{}to gather\textquotedblright{}; MW gives \emph{ava-ci} as to gather fruit from a tree, with the participle \emph{avacita}, \textquotedblleft{}gathered\textquotedblright{} (MW).
\end{cousinweb}

\begin{grammarlens}[title={What you've built}]
One more word for getting around.
\end{grammarlens}

\subsection*{Your turn}
\begin{itemize}
\raggedright
  \item \emph{Say it:} \emph{avacinoti}
  \item \emph{Say it:} \emph{avacinoti}, once more
  \item \emph{Say it:} \emph{bālikā vṛkṣāt phalāni avacinoti.} — The girl picks fruit from the tree.
  \item \emph{Recall:} say the Sanskrit for fertile land, then the Sanskrit for reaps a crop, then say \emph{avacinoti} again
\end{itemize}

\begin{tcolorbox}[breakable,colback=teal!4,colframe=teal!35!black,title={Before you move on}]
What does \sk{अवचिनोति} mean? (\textquotedblleft{}Picks fruit; plucks flowers\textquotedblright{}.) Say it once more.
\end{tcolorbox}
\section[(dialogue)]{Looking back: from Bad Weather to Fields on the Way Home}

\label{lesson:SA-R447-review-441-447}



\begin{quote}
Say the Sanskrit for a day of bad weather and rain-laden.
\end{quote}

\subsection*{Your turn}
Say these aloud:

\begin{itemize}
\raggedright
  \item a day of bad weather, rain-laden, too much rain, hoarfrost, veils
  \item sunrise, rises in the sky, grows bright, the horizon, dawns
  \item sunset, sets in the west, dusk, a dark night, by night
  \item the full-moon day, the new-moon night, the bright fortnight, the dark fortnight, gleams
  \item a shooting star, the pole star, an eclipse, by day, passes the time
  \item a plateau, a summit, a thicket, kusha grass, a young shoot
  \item a flooded field, a crop, fertile land, reaps a crop, picks fruit
\end{itemize}

\begin{tcolorbox}[breakable,colback=teal!4,colframe=teal!35!black,title={Before you move on}]
A day of bad weather? (\textbf{\sk{दुर्दिनम्}}, \emph{durdinam}.) Rain-laden? (\textbf{\sk{सजलः}}, \emph{sajalaḥ}.) Too much rain? (\textbf{\sk{अतिवृष्टिः}}, \emph{ativṛṣṭiḥ}.) Hoarfrost? (\textbf{\sk{तुहिनम्}}, \emph{tuhinam}.) Veils? (\textbf{\sk{आवृणोति}}, \emph{āvṛṇoti}.) Sunrise? (\textbf{\sk{सूर्योदयः}}, \emph{sūryodayaḥ}.) Rises in the sky? (\textbf{\sk{उदेति}}, \emph{udeti}.) Grows bright? (\textbf{\sk{प्रकाशते}}, \emph{prakāśate}.) The horizon? (\textbf{\sk{क्षितिजम्}}, \emph{kṣitijam}.) Dawns? (\textbf{\sk{विभाति}}, \emph{vibhāti}.) Sunset? (\textbf{\sk{सूर्यास्तः}}, \emph{sūryāstaḥ}.) Sets in the west? (\textbf{\sk{अस्तं} \sk{गच्छति}}, \emph{astaṃ gacchati}.) Dusk? (\textbf{\sk{गोधूलिः}}, \emph{godhūliḥ}.) A dark night? (\textbf{\sk{तमिस्रा}}, \emph{tamisrā}.) By night? (\textbf{\sk{नक्तम्}}, \emph{naktam}.) The full-moon day? (\textbf{\sk{पूर्णिमा}}, \emph{pūrṇimā}.) The new-moon night? (\textbf{\sk{अमावास्या}}, \emph{amāvāsyā}.) The bright fortnight? (\textbf{\sk{शुक्लपक्षः}}, \emph{śuklapakṣaḥ}.) The dark fortnight? (\textbf{\sk{कृष्णपक्षः}}, \emph{kṛṣṇapakṣaḥ}.) Gleams? (\textbf{\sk{राजते}}, \emph{rājate}.) A shooting star? (\textbf{\sk{उल्का}}, \emph{ulkā}.) The pole star? (\textbf{\sk{ध्रुवः}}, \emph{dhruvaḥ}.) An eclipse? (\textbf{\sk{ग्रहणम्}}, \emph{grahaṇam}.) By day? (\textbf{\sk{दिवा}}, \emph{divā}.) Passes the time? (\textbf{\sk{यापयति}}, \emph{yāpayati}.) A plateau? (\textbf{\sk{अधित्यका}}, \emph{adhityakā}.) A summit? (\textbf{\sk{शिखरम्}}, \emph{śikharam}.) A thicket? (\textbf{\sk{गुल्मः}}, \emph{gulmaḥ}.) Kusha grass? (\textbf{\sk{कुशः}}, \emph{kuśaḥ}.) A young shoot? (\textbf{\sk{पल्लवः}}, \emph{pallavaḥ}.) A flooded field? (\textbf{\sk{केदारः}}, \emph{kedāraḥ}.) A crop? (\textbf{\sk{सस्यम्}}, \emph{sasyam}.) Fertile land? (\textbf{\sk{उर्वरा}}, \emph{urvarā}.) Reaps a crop? (\textbf{\sk{लुनाति}}, \emph{lunāti}.) Picks fruit? (\textbf{\sk{अवचिनोति}}, \emph{avacinoti}.)
\end{tcolorbox}
\section[(dialogue)]{Fields on the Way Home — once more}

\label{lesson:SA-R447-review-final-447}



\begin{quote}
Say the Sanskrit for a flooded field and a crop.
\end{quote}

\subsection*{Your turn}
\begin{itemize}
\raggedright
  \item \emph{Say it:} a flooded field, a crop, fertile land, reaps a crop, picks fruit
\end{itemize}

\begin{tcolorbox}[breakable,colback=teal!4,colframe=teal!35!black,title={Before you move on}]
A flooded field? (\textbf{\sk{केदारः}}, \emph{kedāraḥ}.) A crop? (\textbf{\sk{सस्यम्}}, \emph{sasyam}.) Fertile land? (\textbf{\sk{उर्वरा}}, \emph{urvarā}.) Reaps a crop? (\textbf{\sk{लुनाति}}, \emph{lunāti}.) Picks fruit? (\textbf{\sk{अवचिनोति}}, \emph{avacinoti}.)
\end{tcolorbox}
